% !TEX root = ../nc-cubic.tex
\subsection{二捩れを持つ全対角三様体の帰着}
\label{nds:sec:signs}
\begin{theorem}\label{nds:thm:signs}
前節の五変数で二捩れの交換係数成分が非自明とする。
全 coherent 圏が smooth proper で、残余圏が狭い主予想の
変形連結性を満たすなら、変数の置換と次数付き twist の後、
全残余圏は次の表示を持つ。$x,y,z$ は対ごとに反可換、
$u,v$ は中心変数で、中心三次は
\[
 f=\lambda xyz+x^2L_x(u,v)+y^2L_y(u,v)+z^2L_z(u,v)+g(u,v).
\]
$\lambda\ne0$、各 $L_i$ は非零一次式、$g$ は平方因子を持たない二変数三次である。
この表示で全 coherent 圏が smooth proper なすべての場合に
$(\dim HH_1,\dim HH_0)=(5,2)$ となる（残余圏の値）。
\end{theorem}
\begin{proof}
系\ref{nds:cor:six}の有限群表示を用いる。交換比の radical で割った
重み群 $W=W_2\oplus W_3$ は各素数部分で非退化交代形式を持つ。
可換側の三次 $F$ は同じ支持を持ち $W$-次数零であり、$Y=V(F)$ は滑らかである。
$W_2^*$ の非自明元が負にする座標数を $k$ とする。
その負の固有空間の射影 $\PP^{k-1}$ は奇次数性から $Y$ に含まれる。
ここでは正の座標に関する $5-k$ 個の二次偏導関数だけが残り得る。
$5-k<k$ なら共通零点があり特異になる。よって $k\le2$ である。
重み二以下の元のみからなる二元線型符号は次元二以下である。
実際、単元ベクトル $e_i$ があれば支持が $i$ を避ける符号語の
重みは高々一で、その部分空間は次元一以下となる。なければ非零語は
すべて二元支持で、二つの異なる辺は交わる。その和は同じ三角形の第三辺で、
この三辺すべてと交わる別の辺はない。
非退化性により次元は偶数なのでちょうど二である。
従って二つの座標を独立に反転する型 I と、三座標の偶数個を反転する型 II だけがある。

型 I で反転座標を $x,y$、他を $u,v,w$ とすれば
$F=x^2L_1+y^2L_2+g_3(u,v,w)$ である。
直線 $\PP(x,y)$ 上の勾配により $L_1,L_2$ は独立である。
各 $L_i$ の変数の三成分の重みは対応する $x,y$ の重みに等しい
（$-2=1\bmod3$）。従って全三成分の重みは残る三座標の重みの中にある。
型 II の符号指標は $x,y,z$ で相異なる三つの非零指標、$u,v$ で零である。
このとき符号不変単項式は定理の形で尽くされる。$\lambda=0$ なら
平面 $\PP(x,y,z)$ が含まれて特異点を生じる。各座標点での勾配により
$L_i\ne0$、直線 $\PP(u,v)$ 上での勾配により $g$ は平方因子を持たない。
従ってこの型の三成分の重みは $u,v$ の二つの重みの中にある。

いずれも $W_3$ の生成元は高々三つで、その次元は偶数なので高々二である。
三成分の重みがアフィン直線上で $a+t_i b$ と書ける場合、交代形式による
交換指数は $c(t_j-t_i)$、$c=b_3(a,b)$ である。
重みの線型汎関数 $s_i=-b_3(a,w_i)$ で変数を $\zeta^{s_i}$ 倍して
Zhang twist を取るとこの成分は消える。
$F$ と元の $f$ の単項式は重み零なのでこの自己同型は $f$ を保ち、
その三乗は恒等である。したがって任意の中心三次商と全残余圏の同値を保つ
（定理\ref{nqt:thm:reduction}の自由加群シフトの議論）。
型 II では常にこれが適用できる。残るのは純符号の場合と、
型 I で $u,v,w$ の三成分の重み $A,B,C$ がアフィン独立な場合である。
後者では $L_1,L_2$ の独立性により $x,y$ の三成分を順に $A,B$ と取れる。

$HH_1$ を計算する。可換側の恒等 sector では
五つの線型留数の重みは座標指標に全行列式を掛けたものとなる。
型 I は符号成分だけで不変部分が零である。
型 II の純符号では $u,v$ の二つが残る。
非恒等 sector の固有空間で三次の制限が恒等零なら固定成分は射影空間で
$HH_1$ は零。固有値の三乗が $1$ なら、その空間の外の偏導関数は
そこで消えるので、非空な切断は滑らかな三次となる。
$HH_1$ への寄与は三座標ちょうどの平面楕円曲線だけである。
ここで扱う純符号作用とアフィン独立型では射影スカラーの非自明元はない。
三元集合 $T$ の楕円曲線の留数指標は $\sum_{i\in T}w_i$ である。
ねじれ固定点公式 \cite[Theorem~6]{QFOBaranovskyPetrov} において
これを相殺する sector は、非退化性により一意である。
各素数成分での固有指数は $w_j$ と $\sum_{i\in T}w_i$ の交代対合
（全体の符号規約を除く）である。従って $T$ 上で符号指数が零となり、
その元の一つの固有値の全レベル集合がちょうど $T$ となることを要する。

純符号の型 I では $T$ に $x,y$ が両方あれば両方が負になる。
一方だけなら正固有空間は四座標で、どちらもなければ恒等元だから
寄与は零。混合型 I では三種類 $A,B,C$ が一つずつなら、アフィン独立性により
その和は零でなく、三種類が同じ指数になることはない。
二種類を $2:1$ で使えば和は二種類の差で、ちょうどその二種類が同じ指数となる。
この集合を符号の正固有空間と交差させると、許される $T$ は
$\{x,u,v\},\{x,u,w\},\{y,v,u\},\{y,v,w\}$ の四つである。
従って型 I の $\dim HH_1$ は零または四で、可換三次三様体と変形連結にならない。

純符号の型 II では各非恒等元の固定軌跡は、負の二座標の直線と、
残る三座標での楕円曲線である。その留数指標は当該元を消す唯一の
非自明指標なので、ねじれ指標と一致して各元が $HH_1$ に一次元寄与する。
$HH_0$ では各非恒等元の両成分ともに幾何的作用は自明であり、
非自明なねじれ指標を掛けると寄与は零。恒等 sector は四次元で、
唯一の既約射影表現に対応する例外二項を除いて二次元となる。
これで主張を得る。付属計算器
\texttt{verify\_diagonal\_sign\_sectors.py} は型 I の全 $432$ 個の
アフィン座標枠と純符号二型について同じ計数を行う。
\end{proof}

\begin{proposition}[残る型の変形連結性]\label{nds:prop:sign-path}
定理\ref{nds:thm:signs}の型 II の任意の smooth proper な中心三次商は、
全 coherent 残余圏が実際に滑らかな可換三次三様体の残余圏と
連結な smooth proper 線型圏族で結ばれる。
\end{proposition}
\begin{proof}
$T=\mathbb A^1\setminus V((t^3-1)(t^3+8))$ とする。
$\{y,z\}=-t x^2$ とその巡回式の定める代数を $S_t$ と書く。
三平方 $X=x^2,Y=y^2,Z=z^2$ は中心である。実際、平方の交換子を
関係式で巡回させると自分自身の $t^3$ 倍に戻る。
従って $S_t$ は $X,Y,Z$ を対角成分、非対角成分を
$q_{12}=-tZ/2,q_{13}=-tY/2,q_{23}=-tX/2$ とする三元 Clifford 代数である。
両表示の生成元を対応させれば相互に逆な写像となる。
Clifford 基底により係数環上次数別自由で、中心体積元
$\omega_t=xyz-q_{12}z+q_{13}y-q_{23}x$ は $\omega_0=xyz$ を満たす。
中心変数 $u,v$ を加え、次の中心式で割る。
\[
 f_t=\lambda\omega_t+(1-t/2)\bigl(XL_x+YL_y+ZL_z\bigr)+g(u,v).
\]
$t=0$ は指定した型 II で、$t=2$ では分離和になる。
$t\ne0$ の周囲代数 $S_t$ は実楕円 Sklyanin domain であり、
$2(1-t^3)\omega_t$ はその THCR 商を与える中心元である
（補題\paperref{mixed}{eco:lem:clifford-input}と
命題\paperref{mixed}{ifb:prop:family}の体積元恒等式）。
$t=0$ は反可換三変数量子多項式環で domain である。
各ファイバーで $f_t$ は非零正則であり、各次数の積写像は有限自由
加群間で一定階数を持つ。このため商 $A_T$ は次数別平坦で底変更を保つ。
$t=2$ は $g$ が平方因子を持たず $\lambda\ne0$ なので、
定理\ref{csc:thm:comparison}の楕円一因子と中心二変数の場合である。
この点では全 coherent 残余圏が滑らかな可換三次三様体のものに同値である。

端点を含む相対 smooth proper 開を構成する。
$A_T$ は次数二の中心元 $X,Y,Z,u^2,v^2$ の多項式環上有限であり、
核のある中心準同型でもよい。二次 Veronese の層化から
固定した $P=\PP^4_T$ 上の有限代数 $\mathcal A$ を得る。
$A_T$ は次数一生成なので、その全 tails は
$\mathcal A$ の準連接加群、coherent tails は coherent 加群に対応する。
局所化後の次数零部分は底上平坦である。
$P\times_T P$ 上の有限代数
$\mathcal D=\mathcal A^{\mathrm{op}}\boxtimes_T\mathcal A$ と
加群 $M=\Delta_*\mathcal A$ はともに $T$ 上平坦である。

各端点で $\Perf(P_t,\mathcal A_t)=D^b(\operatorname{qgr} A_t)$ であり、
この圏は smooth proper である。詳しくは $t=2$ では上記分離比較の
式\eqref{csc:eq:coherent}が全線束による生成も与える。
$t=0$ では $W^*$-ねじれ同変層の表示を用いる。
Serre の生成定理と唯一の既約射影表現により各 coherent 層に
$V\otimes\mathcal O_Y(-m)$ の直和から全射がある。
その核に同じ操作を反復する。滑らかな三様体上の coherent 層間の
Ext は一様に有界で、同変 Ext は有限群不変部分なので同じく有界である。
十分長い解消の境界 Ext 類は消え、元の層は有限な線束複体の直和因子となる。
両端で線束は局所射影である。これは中央二次元を反転した環で、
次数 $1,-1$ の積が両順序で $1$ を含むことから従う。
よって両端の perfect 圏の主張を得る。

$M_t$ は上記圏の対角双加群を表すため、両端で
$\mathcal D_t$-perfect である。具体的には
$\mathcal A_t\otimes\mathcal O_{P_t}(j)$ は $j\in\mathbb Z$ で
非有界準連接導来圏を compact 生成する。両側のその外積も積上の
準連接加群を compact 生成する（ample twist での検出と忘却の随伴）。
外積間では Künneth 同型、$\Delta_*\mathcal A_t$ への Hom では
tensor--Hom 随伴が成り立つので、積上の局所 perfect 核は
二つの perfect 圏の Morita tensor 積に対応し、$M_t$ は恒等核に対応する。
これは定理\paperref{mixed}{cmo:thm:relative-smoothness}の有限中心対角の論証である。

各端点のパラメータを局所環で一つの正則元とする。
底上の平坦性と補題\paperref{mixed}{cmo:lem:projective-lift}により、
各端点上の全ての点で $M$ の有限射影次元が持ち上がる。
perfect locus は開であり（補題\paperref{mixed}{cmo:lem:perfect-locus}）、
その補集合の $T$ への像は射影性により閉である。
その像を除いた開 $U\subset T$ は $0,2$ を含み、既約ゆえ連結である。
この開での $M$ の perfect 性は相対 smoothness を与える。
有限代数の局所射影加群間の Hom は coherent で底上平坦であり、
有限複体を通して射影的な底への pushforward は底上 perfect となる。
従って相対 proper 性も成立する。上の ample twist による生成と
底変更から、各ファイバーは対応する有限代数の perfect 圏である。

最後に標準二項 $\mathcal O,\mathcal O(1)$ を相対的に除く。
両者は上の強次数性によって相対 perfect であり、Hom の底変更が成り立つ。
両端で例外対なので、二つの恒等射からの cone と逆順 Hom の
coherent な支持の像をさらに除けば両端を保ったまま相対例外対になる。
その mutation による右直交は $U$ 上 smooth proper であり、
底変更を保つ。この連結な族の二つの端点は主張した全 coherent 残余圏である。
\end{proof}
